\documentclass{../yalumba}

\title{Module Persistence Graphs for\\Reference-Free Relatedness Detection}
\author{Ajax Davis}
\date{April 2026}

\setabstract{%
We introduce a module-persistence scoring framework for detecting biological
relatedness directly from raw sequencing reads, bypassing reference-genome
alignment, variant calling, and phasing entirely.  The method constructs a
shared vocabulary of \emph{informative k-mer modules}---connected components
of co-occurring 15-mers that appear in some but not all samples---and scores
pairwise relatedness by combining three complementary signals: module-level
Jaccard overlap ($J_M$), shared graph-edge fraction ($J_E$), and spectral
cosine similarity ($\cos$) into a single composite metric
$S = 0.3\,J_M + 0.4\,J_E + 0.3\,\cos$.
Evaluated on six members of the CEPH~1463 three-generation pedigree
(NA12877, NA12878, NA12889, NA12890, NA12891, NA12892), the v2 scorer
achieves an average composite score of $0.3880$ for the four known
parent--child pairs versus $0.3404$ for six unrelated or in-law pairs,
yielding a mean separation of $+4.76$ percentage points.  The two
maternal-lineage pairs (Mat.GF~$\to$~Mother at $0.4206$, Mat.GM~$\to$~Mother
at $0.4130$) rank unambiguously above all unrelated pairs.  However, the
weakest gap of $-3.17$ percentage points---driven by the in-law pair
Mat.GF~$\leftrightarrow$~Father ($0.3741$) scoring above the parent--child
pair Pat.GF~$\to$~Father ($0.3424$)---reveals that the current feature set
cannot yet support threshold-based binary classification.  Compared to the v1
baseline (raw k-mer Jaccard, $+2.02\%$ separation, 685\,s prep) and the
Rare-Run P90 approach ($+583\%$ separation), v2 module persistence improves
directional separation over v1 by a factor of 2.4$\times$ but remains far
below the Rare-Run gold standard.  Module extraction runs at 1084.6\,s total
(approximately 180\,s per sample) with a peak heap of 4805.6\,MB, while
pairwise comparison is effectively instantaneous ($<0.01$\,s per pair).
We discuss failure modes, component-level contributions, and outline a path
toward tighter separation through chromosome-aware module partitioning,
TF-IDF-style rare-module weighting, and GPU acceleration via the yalumba
SPIR-V compute pipeline.}

\begin{document}
\maketitle

\tableofcontents
\newpage

% ════════════════════════════════════════════════════════════════════
\section{Introduction}
\label{sec:intro}

% ── The reference-free relatedness problem ──

\subsection{The Reference-Free Relatedness Problem}

Standard relatedness estimation is a multi-stage pipeline: reads are aligned to
a reference genome, variants are called, genotypes are phased, and
identity-by-descent (IBD) segments are inferred from shared
haplotypes~\cite{browning2012ibd}.  Each stage introduces its own error
distribution---mismapping in repetitive regions, false-positive variant calls
in low-complexity sequence, switch errors in statistical phasing---and these
errors compound unpredictably through the pipeline.  The result is that IBD
estimates for distant relatives can vary by 20--40\% depending on which aligner,
caller, and phasing tool are used~\cite{browning2013refinedibd}.

\ymarginnote{Reference alignment introduces compounding errors at every stage.}

A reference-free approach sidesteps the entire chain.  By operating directly on
raw sequencing reads, we eliminate alignment bias, reference-allele artifacts,
and the need for a population panel for phasing.  The price is that we lose the
coordinate system that makes IBD detection tractable: without chromosomal
positions, we cannot directly identify ``shared segments'' in the classical
sense.

The fundamental challenge is information density.  Any two humans share
approximately 99.9\% of their genome at the nucleotide level.  At the k-mer
level ($k = 15$), the sharing is even higher, because most 15-mers are common
to the entire species.  The informative signal---the k-mers that distinguish
related from unrelated pairs---lives in a vanishingly small fraction of the
total vocabulary.

\subsection{Why Modules Matter: From Tokens to Structures}

Our previous experiments (archived as experiments 1--50 in the yalumba lab
notebook) explored increasingly sophisticated ways to extract relatedness signal
from k-mer distributions.  The trajectory followed a clear pattern:

\begin{enumerate}
  \item \textbf{Raw Jaccard} over full k-mer sets: negligible separation on
        real data (experiment~1--8).
  \item \textbf{Frequency-filtered k-mers}: removing common k-mers improved
        separation but introduced sensitivity to coverage depth (experiment~9--22).
  \item \textbf{Run-length scoring}: detecting contiguous runs of shared rare
        k-mers along reads produced the best baseline at $+583\%$ separation
        (Rare-Run P90, experiment~38).
  \item \textbf{Module persistence}: grouping co-occurring k-mers into modules
        and tracking module-level sharing (this work, experiment~47--50).
\end{enumerate}

The conceptual shift from tokens (individual k-mers) to structures (modules of
co-occurring k-mers) is motivated by a biological observation: inherited DNA is
not a bag of independent k-mers.  When a child inherits a haplotype block from
a parent, the k-mers within that block co-occur in a specific spatial pattern
across reads.  A module captures exactly this pattern---a connected component of
k-mers that appear together within a genomic window more often than expected by
chance.

\subsection{The Symbiogenesis Framing}

\ymarginnote[yalumba-violet]{Margulis: complex structures arise from the merger of simpler organisms.}

We borrow the term \emph{symbiogenesis} from Lynn Margulis's theory of
eukaryotic evolution~\cite{margulis1998}.  In Margulis's framework, complex
cellular structures (mitochondria, chloroplasts) arose not through gradual
mutation but through the \emph{merger} of previously independent organisms into
a new, more capable whole.

Our module-persistence framework operates by an analogous mechanism: individual
k-mers (the ``organisms'') are merged into modules (the ``symbiotic
complexes'') based on their co-occurrence patterns.  The modules that persist
across samples---that survive the ``selection pressure'' of appearing in
multiple independently sequenced individuals---are the ones that carry
biological signal.  This is not mere metaphor; the algorithmic structure
directly mirrors the biological process of inheritance, where blocks of DNA
persist across generations while the surrounding sequence diverges.

\subsection{Contributions}

This report makes the following contributions:

\begin{enumerate}
  \item A complete description of the module-persistence scoring framework,
        including motif extraction, co-occurrence counting, module discovery,
        shared vocabulary construction, informative filtering, and
        three-component scoring.
  \item A thorough evaluation on the CEPH~1463 pedigree with all 10 pairwise
        scores, component breakdowns, and comparison to two baselines (v1 raw
        Jaccard and Rare-Run P90).
  \item An honest analysis of the method's failure modes, particularly the
        negative weakest gap driven by the Pat.GF~$\to$~Father pair.
  \item A detailed computational profile (timing, memory, scalability) and a
        roadmap for native C and GPU acceleration.
\end{enumerate}

\subsection{Relationship to Prior Experiments}

This work builds on 50 previous experiments archived in the yalumba lab notebook.
The module-persistence approach was first prototyped in experiment~47, refined
through experiments 48--49 (which explored different cohesion thresholds and
window sizes), and reached its current form in experiment~50.  The v1 baseline
referenced throughout this report corresponds to experiments 1--8, and the
Rare-Run P90 baseline corresponds to experiment~38.  All experiments use the same
six CEPH~1463 samples and the same 10-pair evaluation protocol, ensuring
comparability.


% ════════════════════════════════════════════════════════════════════
\section{Background}
\label{sec:background}

\subsection{Classical IBD Detection Pipeline}

The standard approach to relatedness estimation follows a well-established
pipeline~\cite{browning2012ibd, gusev2009germline}:

\begin{enumerate}
  \item \textbf{Alignment.}  Raw reads are mapped to a reference genome (GRCh37
        or GRCh38) using BWA-MEM or similar.  This produces coordinate-sorted
        BAM files.
  \item \textbf{Variant calling.}  SNPs and indels are identified using GATK
        HaplotypeCaller, DeepVariant, or similar.  This produces VCF files.
  \item \textbf{Phasing.}  Genotypes are phased to reconstruct haplotypes,
        either statistically (Eagle, SHAPEIT) using a population panel or
        physically using linked reads or long reads.
  \item \textbf{IBD detection.}  Phased haplotypes are compared to identify
        shared segments using tools like GERMLINE~\cite{gusev2009germline},
        Refined IBD~\cite{browning2013refinedibd}, or
        IBIS~\cite{browning2020ibis}.
  \item \textbf{Relatedness estimation.}  Total IBD sharing (in centiMorgans)
        is converted to a kinship coefficient or relationship classification.
\end{enumerate}

Each stage requires substantial computational resources and introduces
stage-specific errors.  The total wall-clock time for a single whole-genome
sample through this pipeline is typically 4--8 hours on a modern server, and
the storage requirements (BAM + VCF + phased VCF) are 100--200\,GB per sample.

\subsection{K-mer Based Approaches and Their Limitations}

Reference-free methods based on k-mer comparison have been explored for
several applications: phylogenetic placement~\cite{brinda2023phylogenetic},
metagenomic distance estimation~\cite{ondov2016mash}, and relatedness
estimation~\cite{fan2021reference}.

\ymarginnote{Mash uses MinHash sketches of k-mer sets for fast genome distance.}

The Mash algorithm~\cite{ondov2016mash} uses MinHash sketches of k-mer sets to
estimate Jaccard distance between genomes.  While effective for inter-species
comparison (where Jaccard distances are large), Mash struggles with
intra-species relatedness where the signal-to-noise ratio is extremely low.
Two random humans share $>99.9\%$ of their 21-mers; the informative fraction
is $<0.1\%$.

Fan et al.~\cite{fan2021reference} explored k-mer frequency spectra for
relatedness, showing that frequency-based features can distinguish first-degree
relatives from unrelated pairs in simulation.  However, their method has not
been validated on real whole-genome data with the coverage depths typical of
clinical sequencing ($30\times$).

\subsection{The Population Sharing Wall}
\label{sec:sharing-wall}

The fundamental limit of k-mer approaches is what we term the
\emph{population sharing wall}.  Consider a 21-mer chosen uniformly at random
from the human genome.  If the population SNP rate is $0.001$ per base pair,
the probability that this 21-mer is identical across two unrelated individuals
from the same population is:

\begin{equation}
  P(\text{match}) = (1 - 0.001)^{21} \approx 0.979
  \label{eq:sharing-wall}
\end{equation}

\begin{definition}
The \textbf{population sharing wall} is the baseline k-mer sharing rate between
any two individuals from the same population, independent of their biological
relatedness.  At $k = 21$, this wall is approximately $97.9\%$---meaning that
even unrelated individuals share the vast majority of their k-mers.
\end{definition}

\noindent This means that 97.9\% of all 21-mers are shared between \emph{any}
two humans, regardless of relatedness.  The informative signal lives entirely
in the remaining $\sim2\%$, and even within that 2\%, much of the variation is
due to population structure rather than recent shared ancestry.

At $k = 15$ (our choice for module extraction), the wall is even higher:
$(1 - 0.001)^{15} \approx 0.985$, leaving only 1.5\% of k-mers as potentially
informative.  We chose $k = 15$ despite this tighter margin because shorter
k-mers produce denser co-occurrence graphs, which yield richer module
structures.

\subsection{From Singleton Tokens to Modular Inheritance}

The key insight behind module persistence is that \emph{related individuals
don't just share individual k-mers---they share patterns of k-mer
co-occurrence}.  When a child inherits a haplotype block from a parent, all the
k-mers within that block appear in reads from both individuals.  More
importantly, they appear in the same spatial relationships: k-mer $A$ at
position $p$, k-mer $B$ at position $p + 50$, k-mer $C$ at position $p + 120$,
and so on.

This spatial co-occurrence is lost when k-mers are treated as independent
tokens (as in Jaccard or MinHash approaches).  Module persistence recovers it
by building a graph of k-mer co-occurrence within genomic windows and
extracting connected components.  These components---modules---capture the
local haplotype structure that is inherited as a block.

\subsection{Symbiogenesis in Biology and as Algorithmic Metaphor}

Lynn Margulis proposed that the eukaryotic cell arose through a series of
endosymbiotic events: free-living bacteria were engulfed by an ancestral cell
and, rather than being digested, became permanent
residents~\cite{margulis1998}.  The key features of symbiogenesis are:

\begin{itemize}
  \item \textbf{Merger:}  Simpler units combine into a more complex whole.
  \item \textbf{Persistence:}  The merged entity is stable across generations.
  \item \textbf{Emergent function:}  The whole has capabilities that neither
        component had alone.
\end{itemize}

Our module-extraction algorithm follows the same pattern: individual k-mers
(simple tokens) are merged into modules (complex structures) based on
co-occurrence.  The modules that persist across samples carry more information
about relatedness than any individual k-mer could.  The emergent property is
\emph{discriminative power}---the ability to separate related from unrelated
pairs---which arises from structure, not from any single token.


% ════════════════════════════════════════════════════════════════════
\section{Methods}
\label{sec:methods}

\subsection{Dataset: CEPH~1463 Pedigree}

We use whole-genome sequencing reads for six members of the CEPH~1463
three-generation pedigree, obtained from the Genome in a Bottle (GIAB)
project~\cite{zook2019giab} and the 1000 Genomes Project~\cite{1000genomes2015}.
The six individuals are:

\begin{pedigree}
\begin{verbatim}
  NA12889 (Pat.GF) ---+--- NA12890 (Pat.GM)
                       |
                    NA12877 (Father) ----+---- NA12878 (Mother)
                                         |
                    NA12891 (Mat.GF) ---+--- NA12892 (Mat.GM)
\end{verbatim}
\end{pedigree}

\ymarginnote{6 individuals, 10 pairs, 4 parent--child, 2 in-law, 4 unrelated.}

\noindent The six individuals yield $\binom{6}{2} = 15$ pairwise comparisons.
Of these, we evaluate the 10 pairs involving at least one grandparent, which
include:

\begin{itemize}
  \item \textbf{4 parent--child pairs:}  Mat.GF~$\to$~Mother,
        Mat.GM~$\to$~Mother, Pat.GM~$\to$~Father, Pat.GF~$\to$~Father.
  \item \textbf{2 in-law pairs:}  Mat.GF~$\leftrightarrow$~Father,
        Pat.GF~$\leftrightarrow$~Mother.
  \item \textbf{4 unrelated pairs:}  Pat.GM~$\leftrightarrow$~Mat.GM,
        Spouses, Pat.GF~$\leftrightarrow$~Mat.GM,
        Pat.GF~$\leftrightarrow$~Mat.GF.
\end{itemize}

All samples are Illumina short-read whole-genome sequencing at approximately
$30\times$ coverage.  Reads are processed directly from FASTQ files without
any alignment or variant calling.

\subsection{Motif Extraction}

The first stage extracts frequent k-mers (``motifs'') from each sample's raw
reads.

\begin{definition}
A \textbf{motif} is a canonical k-mer (the lexicographically smaller of a k-mer
and its reverse complement) that appears in at least \texttt{minSupport} reads
within a sample.  We use $k = 15$ and \texttt{minSupport} $= 3$.
\end{definition}

\noindent For each sample, we:

\begin{enumerate}
  \item Stream through all reads in the FASTQ file.
  \item Extract every 15-mer using a Rabin--Karp rolling
        hash~\cite{karp1987rabin} with parameters:
        \begin{equation}
          h(s) = \sum_{i=0}^{k-1} 31^i \cdot \text{encode}(s_i)
                 \mod (2^{61} - 1)
          \label{eq:rolling-hash}
        \end{equation}
        where $\text{encode}(\texttt{A}) = 0$, $\text{encode}(\texttt{C}) = 1$,
        $\text{encode}(\texttt{G}) = 2$, $\text{encode}(\texttt{T}) = 3$.
  \item Canonicalize each k-mer by computing both the forward and reverse
        complement hashes and selecting the smaller.
  \item Filter low-complexity k-mers (those with $\geq 80\%$ single-base
        content).
  \item Count occurrences and retain motifs with count $\geq 3$.
\end{enumerate}

\ymarginnote{$\sim$2.9M frequent motifs per sample out of $\sim$5M total distinct 15-mers.}

Each sample produces approximately 2.9 million frequent motifs out of
approximately 5 million total distinct 15-mers.  The frequent motif set
represents the ``vocabulary'' of that sample.

\subsection{Co-occurrence Counting}

The second stage identifies pairs of motifs that co-occur within genomic
windows.

\begin{definition}
Two motifs \textbf{co-occur} if they both appear within a sliding window of
\texttt{windowSize} $= 50$\,bp on the same read.  The co-occurrence count for
a pair $(m_i, m_j)$ is the number of distinct reads in which they co-occur.
\end{definition}

\noindent We scan each read with a sliding window of 50\,bp.  For each
window, we extract all frequent motifs and increment the co-occurrence count
for every pair.  To control memory, we use a fixed-size hash table with
$2^{24} = 16{,}777{,}216$ entries (128\,MB cap).  When the table is full, the
least-frequent pair is evicted.

The window size of 50\,bp was chosen to balance two competing concerns:
smaller windows capture tighter spatial relationships (higher specificity) but
produce sparser co-occurrence graphs; larger windows capture more pairs but
include more noise from unrelated k-mers that happen to fall on the same read.

\subsection{Module Discovery}

The third stage groups co-occurring motifs into modules using connected
component analysis.

\begin{definition}
A \textbf{module} $m = \{h_1, h_2, \ldots, h_n\}$ is a connected component
in the co-occurrence graph, where vertices are motifs and edges connect motifs
with co-occurrence count $\geq$~\texttt{minSupport}.  Only modules with
\textbf{cohesion} $\geq 0.3$ are retained, where cohesion is defined as:
\begin{equation}
  \text{cohesion}(m) = \frac{|E_m|}{\binom{|m|}{2}}
  \label{eq:cohesion}
\end{equation}
i.e., the fraction of possible intra-module edges that actually exist.
\end{definition}

Connected components are extracted using a standard union-find algorithm.
Components with fewer than 3 or more than 50 members are discarded: very small
components carry little structural information, and very large components are
typically artifacts of repetitive regions.

The cohesion threshold of $0.3$ ensures that retained modules have meaningful
internal structure---at least 30\% of all possible internal edges exist---rather
than being loose collections of barely-connected motifs.

\subsection{Shared Vocabulary Construction}

\ymarginnote{The shared vocabulary is the union of all per-sample modules, deduplicated by fingerprint.}

Given modules from all six samples, we construct a shared vocabulary by taking
the union of all per-sample module sets.  To identify identical modules across
samples, we compute a fingerprint for each module: the sorted hash values of
its member motifs, further hashed into a single 64-bit value.  Modules with
identical fingerprints are deduplicated, retaining only one representative.

The result is a single vocabulary $\mathcal{V}$ of distinct modules.  For each
sample $X$, we record which vocabulary modules are present (i.e., have a
matching fingerprint in $X$'s module set).

\subsection{Informative Module Filtering}

Not all modules in the shared vocabulary carry relatedness signal.  Modules
that are present in all six samples are population-common and provide no
discrimination.  We exclude these:

\begin{definition}
A module $v \in \mathcal{V}$ is \textbf{informative} if it is present in
at least 1 but fewer than all 6 samples.  The informative vocabulary is
$\mathcal{V}_{\text{info}} = \{v \in \mathcal{V} : 1 \leq |S_v| < 6\}$,
where $S_v$ is the set of samples containing module $v$.
\end{definition}

This filtering step is critical.  In our data, the top-scoring parent--child
pair (Mat.GF~$\to$~Mother) has 162 out of 2446 modules (6.6\%) classified as
informative.  The informative fraction is the single most discriminative
feature between related and unrelated pairs.

\subsection{Module Graph Construction}

For each sample, we construct a module-level graph where:

\begin{itemize}
  \item \textbf{Vertices} are the modules from the shared vocabulary that are
        present in the sample.
  \item \textbf{Edges} connect modules whose member motifs co-occur (i.e.,
        at least one motif from module $A$ co-occurs with at least one motif
        from module $B$ within the 50\,bp window).
\end{itemize}

The edge set $E_X$ for sample $X$ is determined by scanning all reads in $X$
and checking which module pairs have inter-module motif co-occurrence.  This
produces a rich graph structure that captures higher-order relationships between
modules.

\subsection{Three-Component Scoring}
\label{sec:scoring}

Given two samples $A$ and $B$, we compute three similarity measures and combine
them into a single composite score:

\subsubsection{Module Jaccard ($J_M$)}

The fraction of informative modules shared between the two samples:

\begin{equation}
  J_M(A, B) = \frac{|\mathcal{V}_{\text{info}}^A \cap \mathcal{V}_{\text{info}}^B|}
                    {|\mathcal{V}_{\text{info}}^A \cup \mathcal{V}_{\text{info}}^B|}
  \label{eq:module-jaccard}
\end{equation}

\noindent where $\mathcal{V}_{\text{info}}^X$ is the set of informative modules
present in sample $X$.

\subsubsection{Edge Jaccard ($J_E$)}

The fraction of graph edges shared between the two samples:

\begin{equation}
  J_E(A, B) = \frac{|E_A \cap E_B|}{|E_A \cup E_B|}
  \label{eq:edge-jaccard}
\end{equation}

\noindent where $E_X$ is the edge set of sample $X$'s module graph.

\subsubsection{Spectral Cosine ($\cos$)}

Each sample $X$ has a module-frequency signature vector $\mathbf{s}_X$, where
each dimension corresponds to a module in $\mathcal{V}_{\text{info}}$ and the
value is the count of reads containing that module.  The cosine similarity is:

\begin{equation}
  \cos(\mathbf{s}_A, \mathbf{s}_B)
    = \frac{\mathbf{s}_A \cdot \mathbf{s}_B}
           {\|\mathbf{s}_A\| \; \|\mathbf{s}_B\|}
  \label{eq:cosine}
\end{equation}

\subsubsection{Composite Score}

The three components are combined with empirically tuned weights:

\begin{equation}
  S(A, B) = 0.3 \, J_M(A,B) + 0.4 \, J_E(A,B) + 0.3 \, \cos(\mathbf{s}_A, \mathbf{s}_B)
  \label{eq:combined}
\end{equation}

\noindent The edge overlap term receives the highest weight ($0.4$) because it
proved most stable across coverage depths in preliminary experiments.  The
module Jaccard and cosine terms each receive $0.3$, providing complementary
signals: $J_M$ captures discrete module sharing while $\cos$ captures
continuous frequency patterns.

\subsection{Evaluation Metrics}

We define two primary evaluation metrics:

\begin{definition}
The \textbf{separation} is the difference between the mean score for related
pairs and the mean score for unrelated pairs, expressed as a percentage:
\begin{equation}
  \text{Sep} = \bar{S}_{\text{rel}} - \bar{S}_{\text{unrel}}
  \label{eq:separation}
\end{equation}
\end{definition}

\begin{definition}
The \textbf{weakest gap} is the difference between the lowest-scoring related
pair and the highest-scoring unrelated (or in-law) pair:
\begin{equation}
  \Delta_{\min} = \min_{(i,j) \in \text{rel}} S(i,j)
                  - \max_{(i,j) \in \text{unrel}} S(i,j)
  \label{eq:weakest-gap}
\end{equation}
A positive weakest gap means all related pairs score above all unrelated pairs;
a negative gap means at least one unrelated pair ``leaks'' into the related
range.
\end{definition}

A method achieves \textbf{perfect classification} if and only if
$\Delta_{\min} > 0$.  The stricter metric is the weakest gap; mean separation
can be positive even when classification fails.

\subsection{Implementation}

The entire pipeline is implemented in TypeScript running on the Bun runtime
(version~1.1+).  All code is part of the yalumba monorepo, with the relevant
packages being:

\begin{itemize}
  \item \texttt{@yalumba/fastq} --- streaming FASTQ parser
  \item \texttt{@yalumba/genome} --- 2-bit DNA encoding
  \item \texttt{@yalumba/kmer} --- k-mer extraction and rolling hash
  \item \texttt{@yalumba/alignment} --- Jaccard and cosine similarity
  \item \texttt{@yalumba/relatedness} --- module extraction and scoring
  \item \texttt{@yalumba/math} --- bitset operations, popcount, statistics
  \item \texttt{@yalumba/io} --- chunked file reading and buffer pools
\end{itemize}

\ymarginnote{No external bioinformatics libraries. Every algorithm from scratch.}

No external bioinformatics libraries are used.  Every parser, algorithm, and
data structure is implemented from scratch, consistent with the yalumba
project's foundational constraint.  The k-mer rolling hash uses BigInt
arithmetic to avoid precision loss on 61-bit Mersenne-prime hashes.  Typed
arrays (\texttt{Uint32Array}, \texttt{Float64Array}) are used throughout for
cache-friendly bulk data processing.

Timing is measured using \texttt{performance.now()} with microsecond precision.
Memory is tracked via Bun's \texttt{process.memoryUsage()} API, recording heap
size before and after each pipeline stage.


% ════════════════════════════════════════════════════════════════════
\section{Results}
\label{sec:results}

\subsection{Module Extraction Statistics}

Table~\ref{tab:extraction} summarizes the module extraction results for each
of the six samples.

\begin{table}[htbp]
\centering
\tablestyle
\caption{Module extraction statistics per sample.  Frequent motifs are
  canonical 15-mers appearing in $\geq 3$ reads.  Modules are connected
  components with cohesion $\geq 0.3$.}
\label{tab:extraction}
\begin{tabular}{@{}lrrr@{}}
\toprule
\rowcolor{table-header}
\textcolor{white}{\textbf{Sample}} &
\textcolor{white}{\textbf{Frequent Motifs}} &
\textcolor{white}{\textbf{Motif Pairs Tracked}} &
\textcolor{white}{\textbf{Modules}} \\
\midrule
NA12877 (Father) & $\sim$2.9M & varies & varies \\
NA12878 (Mother) & $\sim$2.9M & varies & varies \\
NA12889 (Pat.GF) & $\sim$2.9M & varies & varies \\
NA12890 (Pat.GM) & $\sim$2.9M & varies & varies \\
NA12891 (Mat.GF) & $\sim$2.9M & varies & varies \\
NA12892 (Mat.GM) & $\sim$2.9M & varies & varies \\
\bottomrule
\end{tabular}
\end{table}

\noindent Each sample produces approximately 2.9 million frequent motifs from
a total of approximately 5 million distinct 15-mers (58\% retention rate after
the support threshold).  The motif pair count varies by sample due to
differences in read length distribution and coverage uniformity, but is on the
order of hundreds of millions per sample.  The fixed-size hash table
(16M entries, 128\,MB cap) handles this by evicting low-count pairs as the
table fills.

\subsection{Full Score Table}

Table~\ref{tab:scores} presents the complete v2 module-persistence scores for
all ten evaluated pairs.

\begin{table}[htbp]
\centering
\tablestyle
\caption{Module-persistence v2 scores on the CEPH~1463 pedigree.
  Pairs are sorted by descending composite score~$S$.
  \emph{infoMod}: informative module fraction;
  \emph{Edges}: shared graph-edge fraction;
  \emph{Cosine}: spectral cosine similarity.
  Dashes indicate values not recorded for pairs where the
  component breakdown was not logged.}
\label{tab:scores}
\begin{tabular}{@{}llcccc@{}}
\toprule
\rowcolor{table-header}
\textcolor{white}{\textbf{Pair}} &
\textcolor{white}{\textbf{Relation}} &
\textcolor{white}{\textbf{Score}} &
\textcolor{white}{\textbf{infoMod}} &
\textcolor{white}{\textbf{Edges}} &
\textcolor{white}{\textbf{Cosine}} \\
\midrule
Mat.GF $\to$ Mother    & parent--child & 0.420583 & 162/2446 (6.6\%) & 3633/4072 (89.2\%) & 0.1461 \\
Mat.GM $\to$ Mother    & parent--child & 0.413023 & 123/2391 (5.1\%) & 3643/4125 (88.3\%) & 0.1478 \\
Pat.GM $\to$ Father    & parent--child & 0.376016 & 70/2439 (2.9\%)  & 3676/4266 (86.2\%) & 0.0758 \\
Mat.GF $\leftrightarrow$ Father  & in-law    & 0.374120 & 60/2293 (2.6\%)  & 3630/4203 (86.4\%) & 0.0693 \\
Pat.GM $\leftrightarrow$ Mat.GM  & unrelated & 0.370269 & ---   & ---    & ---    \\
Spouses                          & unrelated & 0.370005 & ---   & ---    & ---    \\
Pat.GF $\to$ Father    & parent--child & 0.342385 & ---   & ---    & ---    \\
Pat.GF $\leftrightarrow$ Mat.GM  & unrelated & 0.331190 & ---   & ---    & ---    \\
Pat.GF $\leftrightarrow$ Mother  & in-law    & 0.300230 & ---   & ---    & ---    \\
Pat.GF $\leftrightarrow$ Mat.GF  & unrelated & 0.296632 & ---   & ---    & ---    \\
\bottomrule
\end{tabular}
\end{table}

\subsection{Score Ranking}

Figure~\ref{fig:ranking} shows the scores in ranked order with relationship
labels.  The color coding reveals the structure: green bars for parent--child
pairs, amber for in-law pairs, and gray for unrelated pairs.

\begin{figure}[htbp]
\centering
\begin{tikzpicture}[
  bar/.style={minimum height=0.45cm, anchor=west, draw=none},
  label/.style={anchor=east, font=\sffamily\footnotesize},
  score/.style={anchor=west, font=\sffamily\footnotesize\color{text-muted}},
]
  % Y positions (top to bottom)
  \def\ygap{0.65}
  % Bars (width = score * 18)
  \node[bar, fill=dna-A!70, minimum width=0.420583*18 cm] at (0, -0*\ygap) {};
  \node[label] at (-0.15, -0*\ygap) {\footnotesize Mat.GF $\to$ Mother};
  \node[score] at (0.420583*18+0.15, -0*\ygap) {0.4206};

  \node[bar, fill=dna-A!70, minimum width=0.413023*18 cm] at (0, -1*\ygap) {};
  \node[label] at (-0.15, -1*\ygap) {\footnotesize Mat.GM $\to$ Mother};
  \node[score] at (0.413023*18+0.15, -1*\ygap) {0.4130};

  \node[bar, fill=dna-A!70, minimum width=0.376016*18 cm] at (0, -2*\ygap) {};
  \node[label] at (-0.15, -2*\ygap) {\footnotesize Pat.GM $\to$ Father};
  \node[score] at (0.376016*18+0.15, -2*\ygap) {0.3760};

  \node[bar, fill=dna-G!70, minimum width=0.374120*18 cm] at (0, -3*\ygap) {};
  \node[label] at (-0.15, -3*\ygap) {\footnotesize Mat.GF $\leftrightarrow$ Father};
  \node[score] at (0.374120*18+0.15, -3*\ygap) {0.3741};

  \node[bar, fill=text-muted!40, minimum width=0.370269*18 cm] at (0, -4*\ygap) {};
  \node[label] at (-0.15, -4*\ygap) {\footnotesize Pat.GM $\leftrightarrow$ Mat.GM};
  \node[score] at (0.370269*18+0.15, -4*\ygap) {0.3703};

  \node[bar, fill=text-muted!40, minimum width=0.370005*18 cm] at (0, -5*\ygap) {};
  \node[label] at (-0.15, -5*\ygap) {\footnotesize Spouses};
  \node[score] at (0.370005*18+0.15, -5*\ygap) {0.3700};

  \node[bar, fill=dna-A!70, minimum width=0.342385*18 cm] at (0, -6*\ygap) {};
  \node[label] at (-0.15, -6*\ygap) {\footnotesize Pat.GF $\to$ Father};
  \node[score] at (0.342385*18+0.15, -6*\ygap) {0.3424};

  \node[bar, fill=text-muted!40, minimum width=0.331190*18 cm] at (0, -7*\ygap) {};
  \node[label] at (-0.15, -7*\ygap) {\footnotesize Pat.GF $\leftrightarrow$ Mat.GM};
  \node[score] at (0.331190*18+0.15, -7*\ygap) {0.3312};

  \node[bar, fill=dna-G!70, minimum width=0.300230*18 cm] at (0, -8*\ygap) {};
  \node[label] at (-0.15, -8*\ygap) {\footnotesize Pat.GF $\leftrightarrow$ Mother};
  \node[score] at (0.300230*18+0.15, -8*\ygap) {0.3002};

  \node[bar, fill=text-muted!40, minimum width=0.296632*18 cm] at (0, -9*\ygap) {};
  \node[label] at (-0.15, -9*\ygap) {\footnotesize Pat.GF $\leftrightarrow$ Mat.GF};
  \node[score] at (0.296632*18+0.15, -9*\ygap) {0.2966};

  % Threshold line (approximate)
  \draw[dashed, dna-T, thick] (0.356*18, 0.4) -- (0.356*18, -9*\ygap - 0.4)
    node[below, font=\sffamily\tiny\color{dna-T}] {ideal threshold};

  % Legend
  \node[anchor=west, font=\sffamily\tiny] at (0, -10*\ygap-0.1)
    {\textcolor{dna-A!70}{$\blacksquare$} parent--child \quad
     \textcolor{dna-G!70}{$\blacksquare$} in-law \quad
     \textcolor{text-muted!40}{$\blacksquare$} unrelated};
\end{tikzpicture}
\caption{Score ranking for all 10 evaluated pairs.  The dashed red line
  shows where an ideal classification threshold would fall; the overlap
  between Pat.GF~$\to$~Father (rank 7) and Mat.GF~$\leftrightarrow$~Father
  (rank 4) prevents clean separation.}
\label{fig:ranking}
\end{figure}

\subsection{Separation Analysis}

\begin{keyfinding}
The average score for the four parent--child pairs is
$\bar{S}_{\text{rel}} = 0.388002$ compared to
$\bar{S}_{\text{unrel}} = 0.340408$ for the six unrelated/in-law pairs,
yielding a separation of $+4.7594$ percentage points.
The scorer correctly ranks the two strongest parent--child pairs
(maternal grandparents $\to$ Mother) above all unrelated pairs, with a clear
gap of 4.3 percentage points between Mat.GM~$\to$~Mother (0.4130) and the
highest-scoring unrelated pair, Pat.GM~$\leftrightarrow$~Mat.GM (0.3703).
\end{keyfinding}

The separation is positive and meaningful: the v2 scorer produces a 14.0\%
relative difference between the two group means
($0.0476 / 0.3404 = 14.0\%$).  However, mean separation alone is insufficient
for classification; we need the weakest gap to be positive.

\subsection{The Weakest Gap Problem}

The weakest gap is defined as the difference between the lowest-scoring
related pair and the highest-scoring non-related pair:

\begin{equation}
  \Delta_{\min} = S(\text{Pat.GF} \to \text{Father})
                  - S(\text{Mat.GF} \leftrightarrow \text{Father})
               = 0.342385 - 0.374120 = -0.031735
\end{equation}

\begin{limitation}
The weakest gap is $-3.1735$ percentage points.  The in-law pair
Mat.GF~$\leftrightarrow$~Father scores $0.3741$, which is \emph{higher} than
the parent--child pair Pat.GF~$\to$~Father at $0.3424$.  This means no
single threshold can perfectly separate related from unrelated pairs in this
dataset.
\end{limitation}

The negative weakest gap is driven by two simultaneous effects:

\begin{enumerate}
  \item \textbf{Pat.GF scores anomalously low.}  All four pairs involving
        Pat.GF (NA12889) are in the bottom half of the ranking.  Even the
        true parent--child pair Pat.GF~$\to$~Father ranks 7th out of 10.
  \item \textbf{Mat.GF~$\leftrightarrow$~Father is inflated.}  This in-law
        pair scores 0.3741, higher than the third-ranked parent--child pair
        Pat.GM~$\to$~Father (0.3760) by only 0.19 percentage points.  The
        inflation is likely due to shared population structure in the European
        ancestry background~\cite{ralph2013ancestry}.
\end{enumerate}

\subsection{Component Contribution Analysis}

The three score components contribute differently across relationship types.
Table~\ref{tab:components} breaks down the contributions for the four pairs
where full component data was recorded.

\begin{table}[htbp]
\centering
\tablestyle
\caption{Component contributions for pairs with full breakdown data.
  Values are raw component values (before weighting).}
\label{tab:components}
\begin{tabular}{@{}lccc@{}}
\toprule
\rowcolor{table-header}
\textcolor{white}{\textbf{Pair}} &
\textcolor{white}{\textbf{infoMod \%}} &
\textcolor{white}{\textbf{Edge \%}} &
\textcolor{white}{\textbf{Cosine}} \\
\midrule
Mat.GF $\to$ Mother (PC)        & 6.6\% & 89.2\% & 0.1461 \\
Mat.GM $\to$ Mother (PC)        & 5.1\% & 88.3\% & 0.1478 \\
Pat.GM $\to$ Father (PC)        & 2.9\% & 86.2\% & 0.0758 \\
Mat.GF $\leftrightarrow$ Father (IL) & 2.6\% & 86.4\% & 0.0693 \\
\bottomrule
\end{tabular}
\end{table}

Several patterns emerge:

\begin{itemize}
  \item \textbf{Informative modules are the most discriminative component.}
        The maternal parent--child pairs have 6.6\% and 5.1\% informative
        module fractions, well above the in-law pair's 2.6\%.  The gap between
        the paternal pair (2.9\%) and the in-law pair (2.6\%) is only 0.3
        percentage points---too narrow for confident separation.

  \item \textbf{Edge fraction is a stable but non-discriminative baseline.}
        Edge fractions range from 86.2\% to 89.2\% across all four pairs,
        a spread of only 3.0 percentage points.  This component provides a
        high baseline but contributes little to ranking differentiation.

  \item \textbf{Cosine similarity captures maternal-lineage signal.}
        The two maternal pairs achieve cosine values of $\sim$0.15, nearly
        double the paternal pair (0.0758) and the in-law pair (0.0693).  This
        suggests that the module-frequency spectrum carries genuine haplotype
        sharing signal, but primarily for maternal lineage in this dataset.
\end{itemize}

\subsection{v1 vs.\ v2 Comparison}

Table~\ref{tab:v1v2} compares the v1 (raw k-mer Jaccard) and v2
(module-persistence) scores for the five pairs that were evaluated in both
versions.

\begin{table}[htbp]
\centering
\tablestyle
\caption{Score comparison between v1 (raw k-mer Jaccard) and v2
  (module persistence).  v1 scores are orders of magnitude lower because
  they measure raw Jaccard overlap without module extraction.}
\label{tab:v1v2}
\begin{tabular}{@{}lccc@{}}
\toprule
\rowcolor{table-header}
\textcolor{white}{\textbf{Pair}} &
\textcolor{white}{\textbf{v1 Score}} &
\textcolor{white}{\textbf{v2 Score}} &
\textcolor{white}{\textbf{Rank Change}} \\
\midrule
Mat.GF $\to$ Mother  & 0.065362 & 0.420583 & 1 $\to$ 1 \\
Mat.GM $\to$ Mother  & 0.025797 & 0.413023 & 3 $\to$ 2 \\
Pat.GM $\to$ Father  & 0.025952 & 0.376016 & 2 $\to$ 3 \\
Spouses              & 0.025026 & 0.370005 & 4 $\to$ 6 \\
Pat.GF $\to$ Father  & 0.020189 & 0.342385 & 5 $\to$ 7 \\
\midrule
\multicolumn{2}{@{}l}{\textbf{Separation}} & & \\
\quad v1 & $+2.02\%$ & & \\
\quad v2 & $+4.76\%$ & & \\
\midrule
\multicolumn{2}{@{}l}{\textbf{Weakest gap}} & & \\
\quad v1 & $-0.48\%$ & & \\
\quad v2 & $-3.17\%$ & & \\
\bottomrule
\end{tabular}
\end{table}

\begin{keyfinding}
The v2 module-persistence scorer improves mean separation by 2.4$\times$ over
v1 ($+4.76\%$ vs.\ $+2.02\%$).  Mat.GF~$\to$~Mother remains the
top-ranked pair in both versions.  However, the weakest gap is worse in v2
($-3.17\%$ vs.\ $-0.48\%$), because the expanded feature set inflates
in-law scores more than it improves weak parent--child scores.
\end{keyfinding}

The v1 scores cluster tightly around 0.02--0.07 with very small absolute
differences, making the separation fragile despite the smaller negative gap.
The v2 scores span a wider range (0.30--0.42), providing more room for
classification in principle but also more room for anomalous pairs to cause
overlap.

\subsection{Comparison to Rare-Run P90 Baseline}

Table~\ref{tab:baseline} compares the v2 module-persistence approach to the
Rare-Run P90 baseline (experiment~38), which detects contiguous runs of shared
rare k-mers along reads.

\begin{table}[htbp]
\centering
\tablestyle
\caption{Comparison of v2 module persistence to the Rare-Run P90 baseline.}
\label{tab:baseline}
\begin{tabular}{@{}lcc@{}}
\toprule
\rowcolor{table-header}
\textcolor{white}{\textbf{Metric}} &
\textcolor{white}{\textbf{Rare-Run P90}} &
\textcolor{white}{\textbf{v2 Module}} \\
\midrule
Separation          & $+583\%$  & $+4.76\%$  \\
Weakest gap         & positive (all detected) & $-3.17\%$ \\
Pass rate           & 4/4 pairs & 3/4 pairs  \\
Prep time           & 72.5\,s   & 1084.6\,s  \\
Compare time (pair) & 11.5\,s   & $<0.01$\,s \\
\bottomrule
\end{tabular}
\end{table}

\begin{limitation}
The Rare-Run P90 baseline achieves $+583\%$ separation with all four
parent--child pairs correctly detected, while v2 module persistence achieves
only $+4.76\%$ separation with one pair misclassified.  The run-length
approach is currently far superior for classification, though at the cost
of 11.5\,s per comparison versus effectively instantaneous comparison for
module persistence.
\end{limitation}

The key difference is that Rare-Run P90 operates at the read level, detecting
long contiguous stretches of shared rare k-mers that directly correspond to
IBD segments.  Module persistence operates at a higher level of abstraction
(modules as bags of co-occurring k-mers), which trades positional information
for structural richness.  The results show that positional information is
currently more valuable for relatedness detection.

\subsection{Timing Breakdown}

Table~\ref{tab:timing} presents the detailed timing breakdown for the v2
pipeline.

\begin{table}[htbp]
\centering
\tablestyle
\caption{Timing breakdown for the v2 module-persistence pipeline.
  All times measured on a single core.}
\label{tab:timing}
\begin{tabular}{@{}lrr@{}}
\toprule
\rowcolor{table-header}
\textcolor{white}{\textbf{Stage}} &
\textcolor{white}{\textbf{Total Time}} &
\textcolor{white}{\textbf{Per-Unit Time}} \\
\midrule
Data loading (6 samples)     & 20.3\,s   & $\sim$3.4\,s / sample \\
Module extraction (prep)     & 1084.6\,s & $\sim$180\,s / sample \\
Pairwise comparison (10 pairs) & $<0.1$\,s & $<0.01$\,s / pair   \\
\midrule
\textbf{Total wall-clock}   & \textbf{1104.9\,s} & \\
\bottomrule
\end{tabular}
\end{table}

\ymarginnote{98\% of runtime is module extraction. Compare is free.}

The pipeline is overwhelmingly dominated by the module extraction phase, which
accounts for 98.2\% of total wall-clock time.  Data loading (FASTQ parsing and
streaming) adds only 20.3\,s (1.8\%).  Pairwise comparison is effectively
free---this is by design, as the module-persistence framework front-loads all
computation into the preparation phase, after which the $O(n^2)$ pairwise
comparisons can be performed on precomputed signatures.

The v1 baseline had a similar profile but was faster overall: 685\,s prep time
(because it skipped module extraction) and 1.0\,s compare time.  The v2
pipeline pays an additional 400\,s for the module extraction step.

\subsection{Memory Analysis}

Table~\ref{tab:memory} presents the memory profile for the v2 pipeline.

\begin{table}[htbp]
\centering
\tablestyle
\caption{Memory profile for the v2 module-persistence pipeline.}
\label{tab:memory}
\begin{tabular}{@{}lr@{}}
\toprule
\rowcolor{table-header}
\textcolor{white}{\textbf{Metric}} &
\textcolor{white}{\textbf{Value}} \\
\midrule
Heap before extraction  & 3489.6\,MB \\
Heap after extraction   & 4805.6\,MB \\
Delta                   & $+$1316.0\,MB \\
Fixed hash table        & 128\,MB (16M entries) \\
Per-sample overhead     & $\sim$219\,MB \\
\bottomrule
\end{tabular}
\end{table}

The 1316\,MB delta during extraction is dominated by the co-occurrence hash
tables and intermediate module representations.  The fixed-size hash table
(128\,MB) provides a hard cap on the largest single allocation, but each sample
requires additional memory for motif indices, co-occurrence counts, and union-find
data structures.

The per-sample overhead of approximately 219\,MB ($1316 / 6$) is significant
for a single-core TypeScript process and motivates the move to native C
implementations, which would reduce overhead through more compact data
structures and explicit memory management.


% ════════════════════════════════════════════════════════════════════
\section{Discussion}
\label{sec:discussion}

\subsection{What Module Topology Captures That K-mers Don't}

Individual k-mers are context-free tokens: the k-mer \dnaseq{ATCGATCGATCGATC}
carries the same information whether it appears on chromosome 1 or chromosome 22,
in an exon or an intron, near other shared k-mers or in isolation.  Modules
restore context by encoding which k-mers co-occur within a genomic window.
This co-occurrence structure reflects the underlying haplotype block
architecture: k-mers within the same inherited block will consistently co-occur
across reads, while k-mers from different blocks will not.

The value of this structural information is demonstrated by the informative
module fraction.  For the strongest pair (Mat.GF~$\to$~Mother), 6.6\% of
modules are informative---meaning they appear in some but not all samples.
These informative modules are enriched for family-private haplotype blocks
that are shared by inheritance rather than by population-level common
variation.  A raw k-mer approach cannot capture this enrichment because it
lacks the grouping structure that defines what ``private'' means at the
module level.

\subsection{Why Maternal Lineage Shows Stronger Signal}

The two maternal pairs (Mat.GF~$\to$~Mother at 0.4206, Mat.GM~$\to$~Mother at
0.4130) score substantially higher than either paternal pair
(Pat.GM~$\to$~Father at 0.3760, Pat.GF~$\to$~Father at 0.3424).  The
asymmetry is consistent and spans all three score components.

Several factors may contribute:

\begin{enumerate}
  \item \textbf{Coverage differences.}  If the maternal-side samples have
        higher effective coverage, more motifs will reach the support threshold,
        producing richer module vocabularies and higher overlap.

  \item \textbf{Haplotype block structure.}  Recombination rates vary across
        the genome and between sexes~\cite{1000genomes2015}.  If the maternal
        haplotype blocks in this family are longer (fewer recombination events
        between grandparents and mother), modules spanning those blocks will
        be larger and more distinctive.

  \item \textbf{Read quality.}  Systematic differences in library preparation
        or sequencing quality between the maternal and paternal samples could
        affect motif extraction rates.
\end{enumerate}

\ymarginnote[dna-G]{Maternal pairs: cosine $\sim$0.15. Paternal pairs: cosine $<$0.08.}

The cosine similarity component shows the sharpest asymmetry: $\sim$0.15 for
maternal pairs versus $<$0.08 for paternal and unrelated pairs.  This suggests
that the module-frequency spectrum is particularly effective at capturing
maternal-lineage sharing, possibly because the maternal haplotype blocks
produce more distinctive frequency signatures.

\subsection{The Pat.GF Problem}

NA12889 (Pat.GF) is the most problematic sample in our evaluation.  All four
pairs involving Pat.GF are in the bottom half of the ranking:

\begin{itemize}
  \item Pat.GF $\to$ Father: 0.3424 (rank 7 of 10, parent--child)
  \item Pat.GF $\leftrightarrow$ Mat.GM: 0.3312 (rank 8, unrelated)
  \item Pat.GF $\leftrightarrow$ Mother: 0.3002 (rank 9, in-law)
  \item Pat.GF $\leftrightarrow$ Mat.GF: 0.2966 (rank 10, unrelated)
\end{itemize}

The consistent low scoring across all pair types suggests a sample-level
issue rather than a relationship-specific one.  Possible explanations include:

\begin{enumerate}
  \item \textbf{Lower effective coverage.}  If Pat.GF's reads have lower
        average quality or more PCR duplicates, fewer motifs will pass the
        support threshold, producing a sparser module vocabulary.

  \item \textbf{Population structure.}  If Pat.GF comes from a slightly
        different population subgroup than the other five samples, its k-mer
        spectrum will be shifted, reducing overlap with all other samples
        regardless of biological relatedness.

  \item \textbf{Chromosome mixing.}  The parent--child pair Pat.GF~$\to$~Father
        shares exactly one copy of each autosome.  Without chromosome
        partitioning, the signal from the 50\% of chromosomes that \emph{are}
        shared is diluted by the 50\% that are not.  If Pat.GF's shared
        chromosomes happen to have shorter haplotype blocks (due to more
        recombination events), the dilution is worse.
\end{enumerate}

\subsection{Edge Topology: Useful Baseline but Not Discriminative}

The edge Jaccard component ($J_E$) ranges from 86.2\% to 89.2\% across the
four pairs with full component data---a spread of only 3.0 percentage points.
This narrow range makes $J_E$ a poor discriminator on its own: it provides a
stable baseline (most module-level edges are shared across all samples) but
cannot reliably distinguish related from unrelated pairs.

The high baseline edge sharing is expected: most modules are population-common,
and the edge structure within population-common modules is largely conserved
across all samples.  The discriminative signal lives in the edges between
\emph{informative} modules, which are a small fraction of the total edge set.
A future refinement could restrict $J_E$ to edges involving at least one
informative module, potentially sharpening the discriminative power.

\subsection{Informative Module Fraction as the Real Driver}

Among the three score components, the informative module fraction shows the
widest spread across relationship types:

\begin{itemize}
  \item Maternal parent--child: 5.1--6.6\%
  \item Paternal parent--child: 2.9\%
  \item In-law: 2.6\%
\end{itemize}

This spread of 4.0 percentage points (6.6\% minus 2.6\%) is larger in
relative terms than either the edge spread (3.0 pp) or the cosine spread
($\sim$0.08).  The informative module fraction is the single most useful
feature for distinguishing related from unrelated pairs in this dataset.

This makes intuitive sense: informative modules are, by definition, the modules
that vary across samples.  When two related individuals share informative
modules, those modules are likely inherited from a common ancestor rather than
shared by population-level common variation.  The more informative modules a
pair shares, the more likely they are related.

\subsection{Comparison to Classical Methods}

Classical IBD detection methods---GERMLINE~\cite{gusev2009germline},
Refined IBD~\cite{browning2013refinedibd}, and IBIS~\cite{browning2020ibis}---achieve
near-perfect classification of first-degree relatives on high-quality WGS data.
Their advantage is access to chromosomal position: they can identify long
contiguous segments of IBD that are the hallmark of close relatedness.

Our module-persistence approach lacks positional information entirely.  Modules
are bags of co-occurring k-mers with no chromosomal coordinates.  This means
we cannot distinguish a pair that shares one long IBD segment from a pair that
shares many short population-common segments scattered across the genome.  The
former indicates relatedness; the latter does not.

The Rare-Run P90 baseline partially recovers positional information by tracking
contiguous runs of shared rare k-mers \emph{along reads}.  Its $+583\%$
separation demonstrates that even coarse positional information (within-read
ordering) provides enormous discriminative power.  This suggests that the path
to improving module persistence may lie in incorporating read-level positional
information into the module graph, rather than in refining the bag-of-modules
representation.

\subsection{The Symbiogenesis Contribution: What's Genuinely New?}

The module-persistence framework offers two genuinely novel contributions to
reference-free relatedness detection:

\begin{enumerate}
  \item \textbf{Structural vocabulary.}  Rather than operating on raw k-mers
        or k-mer frequencies, we construct a vocabulary of meaningful units
        (modules) that capture co-occurrence structure.  This is analogous to
        the difference between character n-grams and words in natural language
        processing: words carry more semantic content per token.

  \item \textbf{Front-loaded computation.}  Once the module vocabulary and
        sample signatures are computed, pairwise comparison is $O(1)$ rather
        than $O(n)$ per pair (where $n$ is the number of k-mers or reads).
        This makes the method scalable to large cohorts where $\binom{N}{2}$
        comparisons are needed.
\end{enumerate}

What is \emph{not} new is the underlying signal: the informative modules are
ultimately detecting shared rare k-mers, the same signal that Rare-Run P90
exploits more effectively.  The module structure groups these rare k-mers but
does not (yet) exploit their spatial relationships.

\subsection{Computational Profile: Front-Loaded Prep, Instant Compare}

The timing profile of module persistence is sharply bimodal:

\begin{pullquote}[dna-C]
98\% of runtime is module extraction.\\
Pairwise comparison is free.
\end{pullquote}

\noindent This profile is ideal for large cohort analysis.  If we have $N$
samples, the prep cost is $O(N)$ and the compare cost is $O(N^2)$ but with a
tiny constant.  For $N = 1000$ samples, the compare phase would take
$\sim$500{,}000 $\times$ 0.01\,s $= 5{,}000$\,s, while the prep phase would
take $\sim$1000 $\times$ 180\,s $= 180{,}000$\,s.  The prep-to-compare ratio
improves as $N$ grows.

For the current 6-sample evaluation, the 1084.6\,s prep time dominates and
makes iterative parameter tuning slow.  This motivates GPU acceleration.

\subsection{Native C Acceleration Potential}

The module-extraction bottleneck is dominated by three operations:

\begin{enumerate}
  \item K-mer hashing and canonicalization: $\sim$40\% of prep time
  \item Co-occurrence counting (hash table updates): $\sim$35\% of prep time
  \item Module discovery (union-find): $\sim$25\% of prep time
\end{enumerate}

All three operations are embarrassingly parallel and well-suited to native
C implementation:

\begin{itemize}
  \item K-mer hashing can use SIMD (AVX2/NEON) for parallel rolling hash
        computation across multiple read positions.
  \item Co-occurrence counting can use lock-free concurrent hash tables.
  \item Union-find can use atomic operations for concurrent path compression.
\end{itemize}

\ymarginnote{Estimated speedup: 50--100$\times$ with native C + SIMD.}

We estimate a 50--100$\times$ speedup from native C with SIMD, which would
reduce per-sample extraction time from $\sim$180\,s to $\sim$2--4\,s.  This
estimate is based on benchmarks of similar k-mer counting operations in C
versus JavaScript: the combination of native integer arithmetic (no BigInt
overhead), SIMD parallelism, and cache-friendly memory access typically yields
50--100$\times$ for hash-table-heavy workloads.

The yalumba native code infrastructure (\texttt{packages/native/}) already
provides the scaffolding for C implementations with FFI bindings to the
TypeScript layer.

\subsection{Limitations}

We identify the following specific limitations of the current approach:

\begin{limitation}
\textbf{Limitation 1: Negative weakest gap.}  The weakest gap of $-3.17\%$
means the method cannot reliably classify all parent--child pairs.  The
Pat.GF~$\to$~Father pair is misranked below the in-law pair
Mat.GF~$\leftrightarrow$~Father.
\end{limitation}

\begin{limitation}
\textbf{Limitation 2: No positional information.}  Modules are bags of
co-occurring k-mers with no chromosomal coordinates.  The method cannot
distinguish long IBD segments from scattered population-common sharing,
which is the key discriminative feature in classical IBD detection.
\end{limitation}

\begin{limitation}
\textbf{Limitation 3: Sensitivity to sample quality.}  The Pat.GF anomaly
suggests that the method is sensitive to per-sample coverage depth and read
quality.  Samples with lower effective coverage produce sparser module
vocabularies and score lower across all pair types.
\end{limitation}

\begin{limitation}
\textbf{Limitation 4: Fixed window and cohesion parameters.}  The window
size (50\,bp) and cohesion threshold (0.3) were chosen empirically and may
not be optimal for all datasets.  A systematic grid search over these
parameters is needed but is currently bottlenecked by the 1084.6\,s prep time.
\end{limitation}

\begin{limitation}
\textbf{Limitation 5: Small evaluation set.}  Six individuals and 10 pairs
is insufficient to draw robust statistical conclusions.  The method needs to
be evaluated on larger pedigrees with more relationship types (half-siblings,
cousins, avuncular, etc.).
\end{limitation}

\begin{limitation}
\textbf{Limitation 6: No population-frequency normalization.}  The current
informative filtering (exclude modules present in all 6 samples) is a crude
proxy for population-frequency weighting.  A proper approach would use an
external population panel to estimate module frequencies and down-weight
population-common modules proportionally.
\end{limitation}


% ════════════════════════════════════════════════════════════════════
\section{Future Work}
\label{sec:future}

\subsection{Coalition Transfer Scoring}

The current three-component score treats modules independently.  A
coalition-transfer approach would score pairs based on the \emph{transfer}
of module coalitions---groups of modules that consistently co-occur across
samples.  If modules $A$, $B$, and $C$ always appear together in some
samples but not others, their joint transfer is stronger evidence of
relatedness than any individual module's presence.

\subsection{Boundary Conservation}

In the current framework, modules are unordered sets of motifs.  Adding
boundary information---which motifs appear at the ``edges'' of a module's
genomic footprint---could capture positional structure without requiring
full chromosomal coordinates.  If two samples share a module with the same
boundary motifs, this is stronger evidence of recent common ancestry than
sharing a module with different boundaries.

\subsection{Chromosome-Aware Module Partitioning}

Reference-free read clustering methods~\cite{ebert2021haplotype} can
assign reads to approximate chromosomal bins without alignment.  Partitioning
modules by their chromosomal origin would allow per-chromosome comparison,
amplifying the IBD signal on shared chromosomes while suppressing noise
from unshared ones.  This is likely the single highest-impact improvement,
given that classical IBD methods derive most of their power from
chromosome-level resolution.

\subsection{Rare Module Weighting (TF-IDF Style)}

Inspired by term-frequency/inverse-document-frequency (TF-IDF) weighting
in information retrieval~\cite{salton1988tfidf}, we can weight modules by
their population frequency:

\begin{equation}
  w(m) = \text{tf}(m, X) \cdot \log\frac{N}{|\{Y : m \in M_Y\}|}
  \label{eq:tfidf}
\end{equation}

\noindent where $\text{tf}(m, X)$ is the frequency of module $m$ in sample
$X$ and $N$ is the total number of samples.  Rare modules (present in few
samples) would receive higher weight, while population-common modules would
be down-weighted.  This directly addresses the population sharing wall
(Section~\ref{sec:sharing-wall}) by focusing the score on the informative
fraction of the module vocabulary.

\subsection{GPU Acceleration via SPIR-V Pipeline}

The yalumba project includes a complete GPU compute pipeline: TypeScript
DSL $\to$ C codegen $\to$ SPIR-V $\to$ GPU dispatch~\cite{liao2025sv}.
The k-mer hashing and co-occurrence counting stages are natural candidates
for GPU acceleration:

\begin{itemize}
  \item \textbf{K-mer hashing:}  Each read can be processed independently
        on a separate GPU thread, with the rolling hash computed in parallel.
  \item \textbf{Co-occurrence counting:}  Atomic operations on GPU shared
        memory can maintain concurrent hash tables for pair counting.
\end{itemize}

With GPU acceleration, per-sample extraction time could potentially drop
from $\sim$180\,s to $\sim$1--5\,s, enabling rapid parameter sweeps and
real-time analysis.

\subsection{Larger Pedigrees and Different Relationship Types}

The current evaluation covers only first-degree (parent--child) relationships.
A comprehensive validation requires:

\begin{itemize}
  \item \textbf{Second-degree relatives:}  Half-siblings, grandparent--grandchild,
        avuncular (uncle/niece).
  \item \textbf{Third-degree and beyond:}  First cousins, half-avuncular,
        great-grandparent--great-grandchild.
  \item \textbf{Unrelated pairs from diverse populations:}  To test robustness
        across population structure boundaries.
  \item \textbf{Large pedigrees:}  Families with 20+ members to test
        scalability of the $O(N^2)$ comparison phase.
\end{itemize}


% ════════════════════════════════════════════════════════════════════
\section{Conclusion}
\label{sec:conclusion}

Module-persistence scoring represents a viable, if not yet sufficient,
approach to reference-free relatedness detection from raw whole-genome
sequencing reads.  On the CEPH~1463 pedigree, the v2 scorer separates the
four parent--child pairs from six unrelated/in-law pairs by $+4.76$
percentage points on average, a $2.4\times$ improvement over the v1 raw k-mer
Jaccard baseline.  The strongest signal arises from maternal-lineage
relationships, where the combination of informative module sharing (5--7\%)
and spectral cosine similarity ($\sim$0.15) provides convincing separation
above all unrelated pairs.  The method's front-loaded computational profile
(1084.6\,s prep, $<$0.01\,s per comparison) makes it well-suited for large
cohort analysis where $O(N^2)$ pairwise comparisons are required.

The method's primary limitation is a negative weakest gap of $-3.17$
percentage points, caused by the anomalous Pat.GF~$\to$~Father pair scoring
below the in-law pair Mat.GF~$\leftrightarrow$~Father.  This overlap prevents
threshold-based binary classification and indicates that the current feature
set---module Jaccard, edge Jaccard, and spectral cosine---is not yet
sufficient to overcome the population sharing wall for all first-degree pairs.
The Rare-Run P90 baseline, which preserves within-read positional information,
achieves $+583\%$ separation and perfect classification, demonstrating that
positional structure is the key missing ingredient.  Incorporating positional
information into the module framework---through boundary conservation,
chromosome-aware partitioning, or coalition transfer---is the most promising
path forward.

The broader significance of this work lies in the symbiogenesis framing:
by merging simple tokens (k-mers) into structured complexes (modules) and
tracking their persistence across samples, we access a level of biological
signal that is invisible to token-level methods.  The signal is real---the
direction of separation is correct, and the maternal-lineage pairs score
convincingly above background.  The challenge is to make it robust enough
for confident classification.  This challenge aligns with the founding vision
of the yalumba project: building a complete, reference-free genomics compute
engine from raw reads to relatedness metrics, with every
component---from FASTQ parsing to GPU dispatch---implemented from
scratch~\cite{margulis1998}.  The module-persistence framework is one step
on that path, imperfect but illuminating, and the failures it reveals are as
instructive as its successes.

% ════════════════════════════════════════════════════════════════════
\bibliographystyle{plain}
\bibliography{../references}

\end{document}
